% Generated draft: requires mathematical and source review.
\documentclass[11pt]{article}
\usepackage[margin=25mm]{geometry}
\usepackage{fontspec}
\setmainfont{Latin Modern Roman}
\usepackage{amsmath,amssymb,mathtools,unicode-math}
\setmathfont{Latin Modern Math}
\usepackage{enumitem,needspace,xurl,hyperref}
\setlength{\parindent}{0pt}
\setlength{\parskip}{5pt}
\setlength{\emergencystretch}{3em}
\newcommand{\N}{\mathbb{N}}
\newcommand{\Z}{\mathbb{Z}}
\newcommand{\Q}{\mathbb{Q}}
\newcommand{\R}{\mathbb{R}}
\newcommand{\C}{\mathbb{C}}
\newcommand{\F}{\mathbb{F}}
\newcommand{\A}{\mathbb{A}}
\newcommand{\PP}{\mathbb{P}}
\newcommand{\E}{\mathbb{E}}
\newcommand{\eps}{\varepsilon}
\newcommand{\dd}{\,\mathrm{d}}
\begin{document}
\section*{UnsolvedMath ID: 800593}
\textbf{Source problem code:} catalogue entry
\subsection{Definitions and mathematical statement}
For a topological space \(Y\), let \(\Omega(Y)\) denote the family of open \(\omega\)-covers of \(Y\): a family \(\mathcal U\) of open subsets of \(Y\) belongs to \(\Omega(Y)\) if \(Y\notin\mathcal U\) and every finite subset of \(Y\) is contained in some member of \(\mathcal U\). Let \(\Gamma(Y)\) denote the family of infinite open \(\gamma\)-covers of \(Y\): using the nontrivial-cover convention \(Y\notin\mathcal U\), these are the infinite families \(\mathcal U\) of open subsets of \(Y\) such that every point of \(Y\) belongs to all but finitely many members of \(\mathcal U\). The property \(\mathsf S_1(\Omega,\Gamma)\) means that for every sequence \((\mathcal U_n)_{n\geq 1}\) with \(\mathcal U_n\in\Omega(Y)\), there are choices \(U_n\in\mathcal U_n\) such that
\[
\{U_n:n\geq 1\}\in\Gamma(Y).
\]
In particular, the selected family must be infinite.

For infinite subsets \(a,b\subseteq\N\), write \(a\subseteq^* b\) if \(a\setminus b\) is finite. A family \(\mathcal A\subseteq[\N]^{\aleph_0}\) has the strong finite intersection property if every finite intersection of members of \(\mathcal A\) is infinite. An infinite pseudointersection of \(\mathcal A\) is an infinite set \(a\subseteq\N\) such that \(a\subseteq^* b\) for every \(b\in\mathcal A\). The pseudointersection number is
\[
\mathfrak p=\min\bigl\{|\mathcal A|:\mathcal A\subseteq[\N]^{\aleph_0},\ \mathcal A\text{ has the strong finite intersection property and has no infinite pseudointersection}\bigr\}.
\]

The supplied record does not explicitly specify the ambient space or the admissible class of \(X\) and \(D\). Its source title suggests sets of reals, but this does not determine the missing hypotheses. To make the union meaningful, let \(X,D\subseteq T\) in the intended common ambient topological space \(T\), and give \(X\) and \(X\cup D\) their subspace topologies. The following questions are conditional on confirming that intended ambient setting and the cover conventions above; they are not assertions quantified over all topological spaces.

\textbf{Question 1 (Miller, Tsaban).} For all admissible \(X,D\subseteq T\), is the implication
\[
X\models\mathsf S_1(\Omega,\Gamma)\quad\text{and}\quad |D|<\mathfrak p
\quad\Longrightarrow\quad
X\cup D\models\mathsf S_1(\Omega,\Gamma)
\]
true?

\textbf{Question 2.} Assuming Martin's Axiom, is it true for all admissible \(X,D\subseteq T\) that
\[
X\models\mathsf S_1(\Omega,\Gamma)\quad\text{and}\quad |D|=\aleph_1
\quad\Longrightarrow\quad
X\cup D\models\mathsf S_1(\Omega,\Gamma)?
\]
The second question is read as replacing the cardinality restriction in the first question while retaining the assumption on \(X\). The source does not state an additional hypothesis concerning the continuum or specify whether its use of Martin's Axiom includes a separate assumption negating the continuum hypothesis.
\subsection{Short English statement}
The problem asks whether adjoining fewer than \(\mathfrak p\) points preserves \(\mathsf S_1(\Omega,\Gamma)\), and whether adjoining \(\aleph_1\) points preserves it under Martin's Axiom. The intended ambient class of spaces and the precise axiom conventions require confirmation. The supplied registry labels the problem open; that label is not independently verified here.
\subsection{Sources}
\begin{itemize}
\item Supplied local UnsolvedMath registry record, permanent ID \(800593\), source problem code \texttt{catalogue entry}, in the AMR Open Problem Lists collection. The record names Tsaban as proposer and gives the year 2006; its title and background attribute the problem to Miller and Tsaban. The supplied status is \texttt{open}.
\item Tsaban, \emph{Selection Principles and Special Sets of Reals: Open Problems} (2006), Problem 20. Supplied source URL: \url{https://arxiv.org/abs/math/0609601}. The record reports extraction by \texttt{arxiv\_latex} and access on 2026-10-16. These are supplied provenance statements, not a claim of independent source access or verification.
\end{itemize}
\end{document}
